%% ============================================================
%% MODELO DE ATIVIDADE · SIMULAÇÃO COMPUTACIONAL
%%
%% Copie este arquivo para atividades/ (atividade-02.tex, por exemplo),
%% inclua-o no relatorio-aacc.tex com \input{atividades/atividade-02}
%% e escreva. Cada \preencher{…} é uma pendência: troque-o pelo seu
%% texto. Cada \figuraprovisoria também: troque-a pela figura de verdade.
%%
%% Uma atividade é UMA entrega sua, com começo, meio e fim. "Participei
%% do projeto X por dois anos" não é uma atividade: são várias. Divida.
%% ============================================================
\begin{atividade}{Simulação de Z para decidir W}{
  tipo         = simulacao,
  modalidade   = {},   % a chave da modalidade pedida: EXT, PES, COMP, VPC, EVT, PRD, ORG, ENS, GES, CUR, OUT
  alternativa  = {},   % a modalidade alternativa, se o Colegiado não aceitar a primeira
  horas        = {},   % as horas desta atividade, com base nos registros (Parte IV)
  inicio       = {},   % AAAA-MM
  fim          = {},   % AAAA-MM, ou: em andamento
  projeto      = {},
  papel        = {},   % Responsável, Corresponsável ou Colaborador(a)
  supervisor   = {},   % quem confirma: nome e cargo (o supervisor do projeto, o gerente)
  registros    = {},   % os códigos no SOMA, com ponto e vírgula: NBL-14; NRO-PRO-003-7
  competencias = {},   % as competências das DCN que ela desenvolveu (I a VIII): I, III, V
  disciplinas  = {},   % as disciplinas do seu curso ligadas a ela: ELE075 Eletrônica; ELE067 Circuitos
}

\begin{contexto}
\preencher{Que pergunta de projeto a simulação respondeu, e que decisão dependia dela.}
\end{contexto}

\begin{contribuicao}
\preencher{O que você modelou, configurou, rodou e analisou; o que era de outras pessoas (o modelo de partida, os dados).}
\end{contribuicao}

\begin{desenvolvimento}
\fichatecnica{
  ferramenta    = {},   % Software e versão
  fenomeno      = {},   % O que foi simulado
  modelo        = {},   % Modelo (físico e numérico)
  parametros    = {},   % Parâmetros e condições de contorno
  discretizacao = {},   % Malha, passo ou discretização — opcional
  validacao     = {},   % Como foi validada
  arquivos      = {},   % Onde estão os arquivos
}

\preencher{O modelo: as equações e as hipóteses simplificadoras. Os parâmetros e as condições de contorno, com a fonte de cada valor (datasheet, medição, literatura). A discretização e o estudo de convergência. A validação: a comparação com a solução analítica, com uma medição ou com a literatura, e o erro.}

\begin{figure}[htbp]
  \centering
  \figuraprovisoria{O modelo: a geometria com a malha, o circuito simulado ou o diagrama do modelo.}
  % \includegraphics[width=0.9\linewidth]{figuras/modelo.pdf}
  \caption{Modelo usado na simulação de Z.}
  \fonte{Elaborado pelo autor.}
  \label{fig:modelo}
\end{figure}

\begin{figure}[htbp]
  \centering
  \figuraprovisoria{O resultado principal, com eixos, unidades e a comparação com a referência.}
  % \includegraphics[width=0.9\linewidth]{figuras/resultado.pdf}
  \caption{Resultado da simulação de Z comparado com a medição.}
  \fonte{Elaborado pelo autor.}
  \label{fig:resultado}
\end{figure}
\end{desenvolvimento}

\begin{resultados}
\preencher{Os resultados com unidades e incertezas, o erro em relação à validação e a decisão de projeto que saiu deles. O que a simulação não permite concluir.}
\end{resultados}

\begin{evidencias}
% Uma linha por evidência: \evidencia{tipo}{identificador}{o que mostra}{onde conferir}
% \evidencia{Arquivo de projeto}{<projeto>/simulacao/Z.asc}{O modelo e os parâmetros da simulação}{Repositório (acesso a pedido)}
% \evidencia{Arquivo controlado}{NRO-PRO-013-<PN>}{O registro da decisão que a simulação apoiou}{SOMA › Arquivos}
\end{evidencias}

\begin{aprendizados}
\preencher{Que conteúdos (métodos numéricos, fenômenos de transporte, circuitos, controle) você aplicou, e o que a simulação de um problema real ensinou.}
\end{aprendizados}
\end{atividade}
